\chapter{Installation and Execution Guide}
\label{app:install}

\section{Prerequisites}

Running the platform needs an AWS account, a domain whose DNS is hosted on Cloudflare and a MongoDB
Atlas cluster. Building it needs Node.js~22, Docker, Terraform~1.14 and the AWS command-line tools,
with credentials of the AWS account in the default profile and a Cloudflare token exported as
\texttt{CLOUDFLARE\_API\_TOKEN}. The installation of these tools is not described here.

\section{Installation Steps}

The repository is one set of npm workspaces, so one command installs every part, and the build
bundles the API, the functions, the transcoder and the web application.

\begin{lstlisting}[language=bash]
npm ci
npm run build
npm run layer:ffprobe -w @zelos/functions   # ffprobe for the probe layer, needs Docker
docker build -f transcoder/Dockerfile -t zelos-transcoder .
\end{lstlisting}

\section{Running Tests}

\texttt{npm run check} runs the linter with the security rules, the type checks and the 294 unit tests.
The smoke test of the image builds its \name{smoke} stage and runs FFmpeg on generated sources. The
end-to-end suite runs against a deployed environment and signs in with test accounts read from a local
file named by \texttt{E2E\_ACCOUNTS}.

\begin{lstlisting}[language=bash]
npm run check
docker build -f transcoder/Dockerfile --target smoke -t zelos-smoke . && docker run --rm zelos-smoke
E2E_BASE_URL=<web origin> E2E_API_URL=<api origin> E2E_ACCOUNTS=<file> \
  npx playwright test -c e2e/playwright.config.ts
\end{lstlisting}

\section{Production Deployment}

The shared environment, which holds the network and the image registry, is applied once. The image is
then pushed with its commit as tag, and the environment is applied with that tag.

\begin{lstlisting}[language=bash]
terraform -chdir=infra/envs/shared init && terraform -chdir=infra/envs/shared apply
docker tag zelos-transcoder <registry>/zelos-transcoder:<commit> && docker push <registry>/zelos-transcoder:<commit>
terraform -chdir=infra/envs/prod init
terraform -chdir=infra/envs/prod apply -var transcoder_image_tag=<commit>
\end{lstlisting}

After the first apply, the fields of Table~\ref{tab:env-secret} are written into the secret of the API
and the connection string into the secret of the workflow. \texttt{scripts/ensure-indexes.sh prod}
creates the database and search indexes, and runs again after every apply that changes the API code.
The Cloudflare Worker is deployed with \texttt{npx wrangler deploy --env prod} after its target host
and origin secret are set, and \texttt{scripts/deploy-web.sh prod} uploads the web application and
refreshes its distribution. The deployment is checked by playing a video through the CDN host and by
running the end-to-end suite against the production origins.
